\subsection{Évaluation d'incidence de la coordinatisation terminale}
\label{subsec:k-incidencias-residentes}

L'état complet conserve l'orientation, le feuillet, la retenue, la signature de frontière,
la mémoire et le livre dodécaphasique. Sur son sous-réseau compatible, le lecteur signé
est la projection de la coordonnée générée et coïncide avec
\[
 R_{\mathrm{sgn}}=\operatorname{pr}_{U}
 =\Pi_H\circ\operatorname{pr}_{C}
 =\Pi_H\circ\mathcal E_{108}^{90,120}\circ\mathcal R_{12}.
\]
Ces identités opèrent sur l'état enrichi, et non sur la projection
brute des cent huit étapes qui oublie ces champs. Une descente par
cette projection requerrait la constance sur ses fibres, et n'est pas présupposée ici.

Pour expliciter la présentation d'incidence, soit \(\mathscr L\) le livre
orienté de l'état complet. Chaque visite conserve la route, la position APP, le feuillet,
l'instant, la multiplicité et l'incidence de frontière. L'évaluation sur le feuillet
d'addition ou de multiplication détermine \(n_v=r_v+9q_v\), avec \(1\leq r_v\leq9\) ;
le résidu n'est pas une entrée indépendante. Pour
\(E_m=\{9(m-1)+1,\ldots,9m\}\), on définit
\[
 \mathscr V_m=\{v:t(v)\in E_m\},\qquad
 \mathscr T_m=\{\theta:\operatorname{origen}(\theta)\in E_m\},
\]
et les dénombrements
\[
\begin{aligned}
 A_m&=\sum_{v\in\mathscr V_m}w(v)\mathbf1_{\{1,3,5,7\}}(r(v)),\\
 C_m&=-\sum_{\theta\in\mathscr T_m}u(\theta)\sigma(\theta)
     =\sum_{j\geq0}\bigl(\nu^-_{120,m}(j)-\nu^+_{120,m}(j)\bigr),\\
 V_m&=\sum_{v\in\mathscr V_m}w(v)\mathbf1_{\{9\}}(r(v))
                         \epsilon_\partial(v).
\end{aligned}
\]
Ici, \(w\) et \(u\) sont les multiplicités des visites et des traces,
\(\sigma\) conserve son orientation et \(\epsilon_\partial\) vaut \(1\) dans
la couronne et \(-1\) à l'intérieur. Les traces de \(\nu^\pm\) sont énumérées
au moyen de la règle du livre ; la somme entière requiert un support fini ou
une construction qui détermine son sens. Le nombre d'incidences
chronologiques et la longueur réduite \(120(1+3j)\) sont des grandeurs distinctes :
leur relation conserve l'unité de longueur du lecteur.

Avec \([x]_{10}\in\{0,\ldots,9\}\), les quotients de \(x=[x]_{10}+10\lfloor x/10\rfloor\) sont également maintenus. La lecture décimale et l'extraction transversale sont
\[
\begin{aligned}
 z_m&=100[A_m]_{10}+10[C_m]_{10}+[V_m]_{10},\\
 \mathcal E_{\mathrm{cnt}}(\mathscr L)
   &=\bigl((z_m-z_{m+3})_{m=1}^{12},
           (z_m-z_{m+4})_{m=1}^{12},\sum_{m=1}^{12}z_m\bigr).
\end{aligned}
\]
Les indices sont cycliques modulo douze. Ainsi, \(\mathcal E_{\mathrm{cnt}}\)
est la présentation \((D_3z,D_4z,Q(z))\) de l'extracteur, avec la même
convention qu'en \S\ref{subsec:k-transversal}. Ses arguments sont les
dénombrements du livre ; ce ne sont ni \(K\), ni \(U\), ni alpha, ni des coordonnées transversales
attendues. La présentation par dénombrements s'applique lorsque le
livre complet est déployé. Elle n'affirme pas que celui-ci soit récupérable après l'oubli
de l'orientation et de la mémoire. L'évaluation harmonique et l'inversion orbitale
développées ci-après conservent le domaine compatible de la
coordinatisation terminale.
